\documentclass[11pt,a4paper]{article}
\usepackage[utf8]{inputenc}
\usepackage[T1]{fontenc}
\usepackage[margin=1in]{geometry}
\usepackage{hyperref}

\hypersetup{pdftitle={WCAG Guidelines and Standards}}
\title{WCAG Guidelines and Standards}
\author{Manuel Pollazzon}
\date{}

\begin{document}
\maketitle

\section{Part 1: Web Accessibility Principles (The POUR Framework)}

WCAG 2.2 is organized around four foundational principles ---
\emph{Perceivable, Operable, Understandable, Robust} --- each of which
subsumes a set of normative guidelines.

\subsection{Principle 1 --- Perceivable}
Information and user interface components must be presentable to users in ways
they can perceive.
\begin{itemize}
  \item \textbf{1.1 Text Alternatives} --- Provide text alternatives for any
        non-text content.
  \item \textbf{1.2 Time-based Media} --- Provide alternatives for time-based
        media (captions, audio description).
  \item \textbf{1.3 Adaptable} --- Create content that can be presented in
        different ways without losing structure.
  \item \textbf{1.4 Distinguishable} --- Make it easier to see and hear
        content, including separating foreground from background.
\end{itemize}

\subsection{Principle 2 --- Operable}
User interface components and navigation must be operable.
\begin{itemize}
  \item \textbf{2.1 Keyboard Accessible} --- Make all functionality available
        from a keyboard.
  \item \textbf{2.2 Enough Time} --- Provide users enough time to read and use
        content.
  \item \textbf{2.3 Seizures and Physical Reactions} --- Do not design content
        in a way that is known to cause seizures.
  \item \textbf{2.4 Navigable} --- Provide ways to help users navigate, find
        content, and determine where they are.
  \item \textbf{2.5 Input Modalities} --- Make it easier to operate
        functionality through various inputs beyond the keyboard.
\end{itemize}

\subsection{Principle 3 --- Understandable}
Information and the operation of the user interface must be understandable.
\begin{itemize}
  \item \textbf{3.1 Readable} --- Make text content readable and
        understandable.
  \item \textbf{3.2 Predictable} --- Make web pages appear and operate in
        predictable ways.
  \item \textbf{3.3 Input Assistance} --- Help users avoid and correct
        mistakes.
\end{itemize}

\subsection{Principle 4 --- Robust}
Content must be robust enough that it can be interpreted reliably by a wide
variety of user agents, including assistive technologies.
\begin{itemize}
  \item \textbf{4.1 Compatible} --- Maximize compatibility with current and
        future user agents (note: \emph{4.1.1 Parsing} was removed in WCAG
        2.2).
\end{itemize}

\section{Part 2: Compliance Levels and Conformance Framework}

\subsection{The Three Conformance Levels}
Every success criterion in WCAG 2.2 is assigned one of three levels:
\begin{itemize}
  \item \textbf{Level A} --- Minimum baseline. Non-negotiable for any
        accessible system.
  \item \textbf{Level AA} --- Standard target for most legal and procurement
        frameworks (e.g., EU EN 301 549, US Section 508, ADA-aligned policy).
  \item \textbf{Level AAA} --- Enhanced, contextual conformance; not
        recommended as a blanket site-wide target.
\end{itemize}

\subsection{The Five Conformance Requirements}
WCAG 2.2 defines five conformance requirements:
\begin{enumerate}
  \item \textbf{Conformance Level} --- Fully satisfy all requirements of the
        chosen level.
  \item \textbf{Full Pages} --- Conformance applies to full web pages, not
        fragments.
  \item \textbf{Complete Processes} --- Every step in a multi-step process must
        conform.
  \item \textbf{Only Accessibility-Supported Ways} --- Use only
        accessibility-supported technologies.
  \item \textbf{Non-Interference} --- No content may interfere with the
        accessibility of other content.
\end{enumerate}

\section{Sources}
https://www.w3.org/TR/2024/REC-WCAG22-20241212/

\end{document}